\documentclass{article}\usepackage[T1]{fontenc}\begin{document}
MixedMates.All.evaluationDiagram proves the mixed base/payload hexagon for evaluation of the single weighted coproduct span at every finite length, including zero. For any base arrow a and ordinary payload natural transformation T the four edges are the existing SpanPull.All.evaluationBase (at both payloads) and SumChange.evaluationPayload (at both base objects). The source face uses the prescribed cleavage obtained by tagging component cleavages, then extending by payload; the target face is naturality of the packed iterate payload map. The proof also establishes the mixed law for the actual inverse comparison. It retains the existing split endpoints and common universe sizes and asserts no naturality or universal property outside this interface.
\paragraph{Component equation.}
Take $a:A\to D$, $T:Y\to Z$, and an object $x$ of the weighted single-span fibre for $Y$ at $D$.
Write $E$ for the actual evaluation, $b_W,b_I$ for reindexing on the weighted span and on the iterates, and $P_W,P_I$ for the payload changes on those two sides. Superscripts indicate the payload; subscripts indicate the base object. Write $S_x$ for the chosen whole-span mixed cell and $H$ for naturality of the packed iterate payload map. Let $\beta$ denote the forward arrow of the existing base-change isomorphism and $\pi$ the \emph{inverse} arrow of the existing payload-change isomorphism. In left-to-right composition order, the theorem identifies
\[
 E_A^Z(S_x)\,;\,\beta^Z_{P_{W,D}x}\,;\,b_I^Z(\pi_{D,x})
 =
 \pi_{A,b_W^Y x}\,;\,P_{I,A}(\beta^Y_x)\,;\,H_{E_D^Y x}.
\]
This equation is promoted to equality of natural transformations; it is not a statement only about images of objects.


Read eq:telescope in manuscript 458ff36a4798e432dee5947afc402f9903a47f14, lines 929--953. In SpanPull.All.evaluationBase and SumChange.evaluationPayload identify the two stages: inverse comparison from the single span to the tagged fibres, then tagged evaluation. ForwardMixed.All.diagram (\#203) only addresses the opposite comparison; TaggedMixed.All.taggedDiagram (\#202) only addresses the tagged evaluation. Treating the inverse functors as strict inverses loses the unit and counit contributions already detectable in the base-direction mate calculation of \#201.

MixedMates.Frame places chosen equivalences at the four corners and specified invertible comparisons on the four edges. Its two grid cells need not be invertible. The diagonal is the composed comparison along the base-then-payload route. Frame.first compares that route with itself; Frame.second is precisely the forward hexagon re-expressed as comparison to the other route along the source and target mixed cells. MateLaws.comp transfers both triangle diagrams with the same diagonal and the actual inverseSquare mates on each edge. Frame.ofPaths substitutes the common diagonal component between the two equations, giving Frame.backLaw without inverting the source or target mixed cell. Frame.composeLaw uses the same common-diagonal argument with the pasting lemma for the two frames. The middle compatibility is the naturality in the pasting proof, not a separate equality assumption.

In All.presentation, the four forward edges are exactly the corrected SpanPull.All.comparison and the isomorphism from compare\_payload. Reversing that frame uses the old inverseSquare at all four edges. In All.tagged, the second frame is built from the old taggedBase and taggedPayload. Unfold ForwardMixed.Commuting and All.evaluationCube in the final theorem: the source cell is the presented whole-span mixed comparison; the target cell is eqToHom of DecoderCube.natural for ChainPayload.allIterateMap. The base edges are hom components of the actual final evaluationBase and the payload edges are inverse components of the actual final evaluationPayload. No replacement square or independently chosen inverse functor is introduced.

In evidence/mixed-mates-author-fresh, FourCorners.lean uses a one-object category with arrow monoid Z times N. The four equivalences have the same underlying functors and four different units. The two mixed cells have second coordinate one and cannot be inverted. The actual inverse faces have signed coordinates 1, 13, -5, -6, and the tested composite has coordinates (11,1). A naive reversal that ignores the equivalence data miscomputes an individual inverse face. SingleSpan.lean specializes the final diagrams to a constant Boolean base endomorphism with no right inverse, and fibrewise Option, which misses arrows that send none to some. It uses whole-span objects reached from tags zero and two, without claiming to compute an arbitrary chosen inverse object.

The result closes this mixed equation, not the entirety of issue \#105. It does not establish a full bundled pseudonaturality structure, separate identity/composition coherence in the payload variable for all final comparison maps, a wider universe regime, the flattening/multiplication for the telescope, the monad unit comparisons, or algebra recovery. The generic result requires the four corner equivalences and the four invertible edge comparisons; it is not a claim about arbitrary lax changes of presentation.
\end{document}
